
\scanpage{52}
\mspage{8}
On a ainsi la notion de \uline{graphe pondéré} \deleted{\unsure{enrichi}} \inserted{\uline{à repères}} $\Gamma$.

\diagram{\begin{tikzcd}[ampersand replacement=\&,column sep=large,row sep=large]
\vec R_l^+\ \deleted{\unsure{\mathbb D_\infty|D_\infty}}\subset\underset{\sigma_0,\sigma_1}{\vec R}
 \arrow[r]\&\vec A\arrow[d,"o"']\arrow[loop right,"\sigma"]\\
\&S\arrow[r,"\gamma"]\&\mathbb N
\end{tikzcd}}

Groupoïde $(\uline{\mathrm{Gr.pondrep}})$

foncteur
\[
 S_*\mathcal TD_{\mathrm{top}}^\pm\longrightarrow(\mathrm{Grpondrep})
\]
qui sera important sans doute dans l'explicitation de l'opération de ``recollement''.
Il est clair en tous cas que toute description algébrico-combinatoire de
$\unsure{S\mathcal TD_{\mathrm{top}}^\pm}$ devra permettre, pour chaque objet de
\deleted{la catégorie} \inserted{groupoïde-}paradigme, de réconstituer (à isom. uniques près) les \uline{graphes pondérés à repères} correspondants.
\inserted{--- plus précisément, il faut réconstituer le foncteur canonique
$S_*\mathcal TD_{\mathrm{top}}^\pm\longrightarrow(\mathrm{Gr.pond.rep.})$
(et non seulement le foncteur à valeurs dans $(\mathrm{Grbipond})$)}

\uline{Opérations essentielles}

A) Opérations ``d'oubli'' ou de ``\uline{gommage}''.

\textcircled{1} \uline{Gommage des repères} À un objet $(X,\Sigma,R)$, muni en plus d'une partie $R'\subset R$, il associe $(X,\Sigma,R\setminus R')$ (gommage de $R'$).
\mspage{51}
